\chapter{Curve regularization and cubical delta-metricity}
\label{app:ch19-comparison-proofs}

Chapter~19 states four comparison results at the points where they are
needed.  Their proofs are collected here so that the worked curve, toric,
model, and abelian examples remain visible as a single route through that
chapter.  The first two proofs concern regularity on curves; the last two
isolate the obstruction to ordinary delta-metricity of cubical metrics.

The two routes require different geometric data. For a curve function,
start with the residual Picard class when testing ordinary smoothness, and
with the positivity of its Hessian current when constructing a decreasing
smooth approximation. For a cubical metric, first reduce to an ample
symmetric rigidified line bundle, then test the Raynaud residual condition.
Each proof uses the full hypotheses of its cited Chapter~19 statement.

\begin{remark}[Proof and source boundary]
\label{rem:appN-proof-boundary}
The four proof files below preserve the dated source packets used in the
Chapter~19 integration. Their verifier records concern those dated bytes.
The proof introductions and navigation in this appendix form part of the
current editorial text and carry no new independent verification receipt.
Ordinary CLD smoothness, ordinary delta-metricity and harmonic
tropicalization are separate conclusions; the criterion cited in each
section determines which one is being tested.
\end{remark}

\section{Residual Picard obstruction on curves}
\label{sec:appN-residual-picard}

We prove \Cref{prop:p213-residual-picard-smooth-curve}.  The quotient of
Thuillier harmonic germs by logarithms of analytic units is measured at a
type~\(2\) point by the real Picard group of its residual curve.

\begin{proof}
By \Cref{cor:ch19-continuous-zero-current-harmonic}, the function \(h\) is
Thuillier harmonic.  Write \(\mathcal H_C\) for the sheaf of Thuillier
harmonic germs and \(\mathcal F_C\) for the real vector subsheaf generated
by germs \(\log|f|\) with \(f\) an invertible analytic function.  Thuillier's
description, recalled in \Cref{rem:ch08-harmonic-vs-log}, says that
\(\mathcal H_C/\mathcal F_C\) is supported at the type~\(2\) points and
that
\[
       (\mathcal H_C/\mathcal F_C)_x
       \simeq\operatorname{Pic}^0(C_x)\otimes_{\mathbf Z}\mathbf R
\]
there; see \cite[Theorem~2.3.21 and Lemma~2.3.22]{Thu05}.  Under this
isomorphism, the class of the germ \(h_x\) is precisely
\(\operatorname{ob}_x(h)\).

On the other hand, \cite[Remark~15.10]{GJR26} identifies
\(\mathcal F_C\) with the kernel of
\[
 \dprimeddprime:\mathcal A^{0,0}_{\mathrm{CLD}}
                    \longrightarrow\mathcal A^{1,1}_{\mathrm{CLD}},
\]
that is, with the sheaf of ordinary CLD-smooth harmonic functions.  Hence a
type~\(2\) harmonic germ lies in that ordinary smooth kernel exactly when
its residual Picard class vanishes.  Away from type~\(2\), the quotient
sheaf is zero, which proves the assertion for points of types \(1\), \(3\),
and \(4\) without imposing a nonexistent annular coordinate at a type~\(4\)
point.  Since ordinary smoothness and the smooth zero-Hessian equation are
local, \eqref{eq:p213-residual-picard-vanishing} proves the global assertion
on \(W\).

Finally, the two stated sufficient hypotheses for the quotient to vanish
are exactly the special cases of \cite[Theorem~2.3.21]{Thu05} recorded in
\Cref{rem:ch08-harmonic-vs-log}.  Equivalently, the abstract group of
degree-zero line-bundle classes on every \(C_x\) is torsion; no uniform
torsion exponent is required.
\end{proof}


\section{Ordinary smooth monotone regularization on curves}
\label{sec:appN-curve-regularization}

We prove \Cref{thm:p063-current-psh-ordinary-regularization}.  The argument
passes from a finite tree skeleton to moving residual divisors and then
replaces the resulting negative outgoing slopes by convex caps.

\begingroup
\emergencystretch=2em
% Source-aware proof for the exact candidate statement
% p063_curve_cld_monotone_regularization_statement_20260827.tex.
% Date: 27 August 2026.
% This packet is isolated from the live book and has not received a final
% production-Rethlas verdict.

\paragraph{Imported facts and conventions.}
We use Thuillier's potential theory on smooth strictly analytic curves.
In particular, lisse functions factor through finite semistable skeleta by
\cite[Proposition~3.2.4]{Thu05}, lisse subharmonicity is characterized by
the nonnegativity of the graph Laplacian by
\cite[Proposition~3.4.4]{Thu05}, and every subharmonic function is the
pointwise infimum of a decreasing net of lisse subharmonic functions on
each relatively compact open set by \cite[Theorem~3.4.2]{Thu05}.  We use
the local semistable structure and the insertion of finitely many type
$2$ or type $3$ points into a skeleton from
\cite[\S2.2.3 and Lemma~2.2.22(ii)]{Thu05}.

At a type-$2$ point \(v\), define \(\operatorname{div}_v(h)\) to be the
divisor whose coefficient at the residual point \(p_\xi\) is the outgoing
slope \(d_\xi h(v)\).  With this definition the smooth image is the
principal subspace.  If the source uses the globally opposite divisor-sign
convention, both residual maps below are multiplied by \(-1\); the
principal subspace is invariant under this change, and the prescribed
outgoing slopes are unchanged.  The two exact residual-divisor sequences
used below are proved in \cite[Lemma~2.3.22]{Thu05}.  The identification of the
subsheaf generated by local analytic units with the sheaf of ordinary
Chambert--Loir--Ducros smooth harmonic functions is
\cite[Remark~15.10]{GJR26}.  This is the ordinary smooth sheaf, not the
weakly smooth sheaf of harmonic tropicalizations.

\begin{proof}[Proof of \cref{lem:p063-local-tree-affinoid}]
Replace \(C\) first by the smooth projective compactification of the
connected component containing \(x\).  The original analytification is an
open subspace of the compactification, so \(W\) is unchanged for the local
question.  If \(x\) is of type \(2\), \(3\), or \(4\), choose a simple
connected open neighbourhood \(U\subset W\) of \(x\) such that
\[
                         \mathbb H(U):=U\setminus U(K)
\]
is an \(\mathbf R\)-tree.  Such neighbourhoods form a fundamental system by
\cite[Corollary~4.27 and Proposition~5.9]{BPR13Structure}.

Strict affinoid neighbourhoods form a fundamental system on a smooth strict
analytic curve; see the proof of \cite[Corollary~2.1.18]{Thu05}.  Hence
there is a strict affinoid neighbourhood \(A\) of \(x\) with \(A\subset U\).
Let \(Y\) be the connected component of \(A\) containing \(x\).  A strict
affinoid algebra has finitely many connected components, each cut out by an
idempotent.  Thus \(Y\) is a connected strict affinoid domain.  Since
\(x\in A^\circ\) and \(Y\) is open and closed in \(A\), we also have
\(x\in Y^\circ\).

By \cite[Theorem~2.3.8]{Thu05}, after a finite separable extension of \(K\),
the domain \(Y\) admits a strictly semistable formal model.  The extension is
trivial because \(K\) is algebraically closed.  Its skeleton \(\Gamma_0\) is
a finite connected metric graph contained in
\(\mathbb H(Y)\subset\mathbb H(U)\).  An \(\mathbf R\)-tree contains no
embedded circuit, so \(\Gamma_0\) is a finite metric tree.

If \(x\) is of type \(1\), use
\cite[Corollary~4.27(1)]{BPR13Structure} to identify an open neighbourhood
\(U\subset W\) with the standard open unit ball.  In a coordinate \(T\),
choose \(R\in|K^\times|\) with \(|T|_x<R<1\) and put
\[
                         Y=\{|T|\leq R\}.
\]
Then \(Y\) is a connected strict affinoid neighbourhood of \(x\), and its
one-point Gauss skeleton \(\Gamma_0\) is a finite tree.

Let now \(f\) be lisse on a neighbourhood of \(Y\).  By
\cite[Proposition~3.2.4]{Thu05}, after a strictly semistable refinement there
are a retraction \(\tau:Y\to\Gamma\) and a continuous piecewise affine
function \(F\) on \(\Gamma\) such that \(f|_Y=F\circ\tau\).  Take a common
semistable refinement with \(\Gamma_0\) and insert all type-\(2\)
breakpoints; retain every type-\(3\) breakpoint as a metric subdivision
point.  The refined skeleton is still contained in
\(\mathbb H(Y)\subset\mathbb H(U)\); in the type-\(1\) case it lies in the
norm-point tree of the open ball.  Hence it is a finite metric tree.  After
the indicated subdivision, \(F\) is affine on each edge, and every added
interior vertex is of type \(2\) or type \(3\).
\end{proof}

\begin{proof}[Proof of
\cref{lem:p063-tree-rational-strictification}]
If \(\Gamma\) has no edge, take \(G=F\).  Assume henceforth that it has an
edge.  Choose \(b\in\partial\Gamma\), regard \(b\) as the root, and put
\[
                 I=V(\Gamma)\setminus\partial\Gamma.
\]
Define \(\mu:V(\Gamma)\to\mathbf Z\) by
\[
 \mu(v)=
 \begin{cases}
  1,&v\in I,\\
  0,&v\in\partial\Gamma\setminus\{b\},\\
  -\#I,&v=b.
 \end{cases}
\]
Thus \(\sum_v\mu(v)=0\).

For an edge \(e\), let \(\Gamma_e\) be the component of
\(\Gamma\setminus e^\circ\) not containing \(b\), and orient \(e\) from
\(\Gamma_e\) towards \(b\).  Prescribe on \(e\) the slope
\[
                         s_e=\sum_{v\in\Gamma_e}\mu(v)
\]
and normalize the resulting piecewise affine function \(\psi\) by
\(\psi(b)=0\).  The absence of circuits makes these prescriptions
compatible and determines \(\psi\).  If \(v\ne b\), the outgoing slope on
the edge from \(v\) to its parent is
\[
 \mu(v)+\sum_{\substack{e\text{ joining }v\\
                        \text{to a child subtree}}}s_e,
\]
whereas the outgoing slope from \(v\) towards a child subtree is
\(-s_e\).  Therefore
\[
                         \Delta_\Gamma\psi(v)=\mu(v).
\]
The same equality at \(b\) follows from \(\sum_v\mu(v)=0\).  In
particular,
\begin{equation}
\label{eq:p063-unit-interior-laplacian}
                         \Delta_\Gamma\psi(v)=1
                         \qquad(v\in I).
\end{equation}

Choose \(\delta>0\) so small that
\[
                   \delta\lVert\psi\rVert_\infty
                   <\frac{\varepsilon}{2}.
\]
Then \(H=F+\delta\psi\) satisfies
\[
                         \Delta_\Gamma H(v)\geq\delta
                         \qquad(v\in I)
\]
by \eqref{eq:p063-graph-subharmonicity} and
\eqref{eq:p063-unit-interior-laplacian}.  Denote the slope of \(H\) on an
oriented edge \(e\) by \(a_e\), and choose \(q_e\in\mathbf Q\) with
\[
                         |q_e-a_e|<\eta.
\]
There is a unique continuous function \(G\) with \(G(b)=H(b)\) and slope
\(q_e\) on every oriented edge.  If \(d\) is the largest valency of a
vertex and \(L\) is the sum of all edge lengths, then
\[
 |\Delta_\Gamma G(v)-\Delta_\Gamma H(v)|\leq d\eta,
 \qquad
 \lVert G-H\rVert_\infty\leq L\eta.
\]
The second estimate follows by integration along the unique path from
\(b\).  Taking
\[
              d\eta<\frac{\delta}{2},
              \qquad
              L\eta<\frac{\varepsilon}{2}
\]
proves \eqref{eq:p063-graph-rational-error} and
\eqref{eq:p063-graph-strict-mass}.
\end{proof}

\begin{proof}[Proof of \cref{lem:p063-residual-divisor-moving}]
After replacing a denominator-clearing integer by a sufficiently large
multiple, choose \(N\) so that \(NB\) is integral and
\[
                         \deg(NB)\geq2g+1.
\]
Put \(L=\mathcal O_{C_0}(NB)\).  For every \(p\in C_0\), Serre duality
gives
\[
 H^1(C_0,L(-p))
 \simeq
 H^0\bigl(C_0,\omega_{C_0}\otimes L^{-1}(p)\bigr)^\vee.
\]
The line bundle on the right has degree
\[
                   2g-1-\deg L<0,
\]
so \(H^1(C_0,L(-p))=0\).  The exact sequence
\[
 0\longrightarrow L(-p)\longrightarrow L
   \longrightarrow L|_p\longrightarrow0
\]
therefore makes \(H^0(C_0,L)\to L|_p\) surjective.  Its kernel is a
proper linear subspace.

For \(p\in M\), choose a nonzero \(k_0\)-linear functional
\(\lambda_p:H^0(C_0,L)\to k_0\) representing evaluation after a
trivialization of \(L|_p\).  The polynomial
\[
                         \prod_{p\in M}\lambda_p
\]
is nonzero.  A nonzero polynomial on a finite-dimensional vector space
over an infinite field cannot vanish everywhere; this follows by
induction on the dimension from the corresponding assertion in one
variable.  We may choose \(s\in H^0(C_0,L)\) with \(s(p)\ne0\) for every
\(p\in M\).  Its zero divisor \(A=(s)_0\) is effective, is disjoint from
\(M\), and satisfies
\[
                         \mathcal O_{C_0}(A)\simeq L
                         \simeq\mathcal O_{C_0}(NB).
\]
Thus \eqref{eq:p063-moving-divisor} holds.  Moreover,
\[
                         N\left(B-\frac1N A\right)=NB-A
\]
is principal.
\end{proof}

\begin{proof}[Proof of \cref{lem:p063-principal-residual-germ}]
Let \(\mathcal H_v\) be the stalk of Thuillier harmonic germs and let
\(\mathcal F_v\) be the stalk generated over \(\mathbf R\) by germs
\(\log|f|\), with \(f\) an invertible analytic function.  The proof of
\cite[Lemma~2.3.22]{Thu05} gives exact sequences
\begin{align}
 0&\longrightarrow\mathbf R\longrightarrow\mathcal H_v
 \xrightarrow{\operatorname{div}_v}
 \operatorname{Div}(C_v)\otimes_{\mathbf Z}\mathbf R
 \xrightarrow{\deg}\mathbf R\longrightarrow0,
\label{eq:p063-harmonic-residual-exact}\\
 0&\longrightarrow\mathbf R\longrightarrow\mathcal F_v
 \xrightarrow{\operatorname{div}_v}
 \operatorname{Prin}(C_v)\otimes_{\mathbf Z}\mathbf R
 \longrightarrow0.
\label{eq:p063-smooth-residual-exact}
\end{align}
The coefficient at \(p_\xi\) is the outgoing slope \(d_\xi h(v)\).
By \eqref{eq:p063-residual-degree-zero}, the divisor \(D\) lies in the
image of the first residual map.  Its class in
\(\operatorname{Pic}^0(C_v)\otimes_{\mathbf Z}\mathbf R\) vanishes
because it is a rational multiple of a principal divisor.  Comparing
\eqref{eq:p063-harmonic-residual-exact} and
\eqref{eq:p063-smooth-residual-exact} gives a germ in \(\mathcal F_v\)
with residual divisor \(D\); adjusting by a constant gives a germ \(h\)
with \(h(v)=c\).
The construction in the proof of Thuillier's lemma permits a
representative containing short arms in every one of the finitely many
displayed directions.

By \cite[Remark~15.10]{GJR26}, \(\mathcal F\) is exactly the sheaf of
ordinary Chambert--Loir--Ducros smooth harmonic functions on a smooth
boundaryless analytic curve.  Hence \(h\) has all the properties in
\eqref{eq:p063-prescribed-residual-slopes}.
\end{proof}

\begin{proof}[Proof of
\cref{prop:p063-tree-lisse-ordinary-regularization}]
By \cite[Proposition~3.2.4]{Thu05}, after a semistable refinement and a
finite metric subdivision, there are a finite metric tree \(\Gamma\), a
retraction \(\tau:Y\to\Gamma\), and a continuous edgewise affine function
\(F\) such that
\[
                         f=F\circ\tau.
\]
Here no extension of scalars is needed because \(K\) is algebraically
closed.  Include every breakpoint and branch point among the metric
vertices, and let \(\partial\Gamma\) contain the vertices over the Shilov
boundary of \(Y\).  Interior vertices are of type \(2\) or type \(3\).
Subharmonicity and \cite[Proposition~3.4.4]{Thu05} give
\[
                         \Delta_\Gamma F(v)\geq0
                         \qquad
                         \bigl(v\notin\partial\Gamma\bigr).
\]
If \(\Gamma\) has no edge, then \(f\) is constant and there is nothing to
prove.  Otherwise apply
\cref{lem:p063-tree-rational-strictification}, with error
\(\varepsilon/3\), to obtain a rational-slope function \(G\) satisfying
\begin{equation}
\label{eq:p063-tree-graph-error}
 \lVert G\circ\tau-f\rVert_{L^\infty(Y)}
 <\frac{\varepsilon}{3}
\end{equation}
and
\begin{equation}
\label{eq:p063-tree-positive-vertex-mass}
                         \Delta_\Gamma G(v)>0
                         \qquad
                         \bigl(v\notin\partial\Gamma\bigr).
\end{equation}

We first modify \(G\circ\tau\) near each interior type-\(2\) vertex
\(v\).  Choose pairwise disjoint simple star neighbourhoods
\(U_v\Subset Y^\circ\).  Such a neighbourhood contains short arms in all
graph directions, and all but finitely many other tangent directions are
open-disc components contained in \(U_v\).  Let \(E_v\) be the finite set
of exceptional directions.  Enlarge
\[
 M_v=\{p_{v,e}:e\ni v\}\subset C_v(\widetilde K)
\]
by the residual points corresponding to \(E_v\).  This is the finite
forbidden set; in particular it contains every graph direction and every
direction in which a fresh collar is not available.

Define
\[
                         B_v
   =\sum_{e\ni v}d_eG(v)[p_{v,e}]
   \in\operatorname{Div}_{\mathbf Q}(C_v).
\]
By \eqref{eq:p063-tree-positive-vertex-mass},
\[
                         \deg B_v
                         =\Delta_\Gamma G(v)>0.
\]
Apply \cref{lem:p063-residual-divisor-moving} to \(B_v\) and \(M_v\).
There are \(N_v>0\) and an effective divisor
\[
                         A_v=\sum_{p\in S_v}m_{v,p}[p]
\]
such that
\[
 A_v\sim N_vB_v,
 \qquad
 S_v\cap M_v=\varnothing.
\]
Thus
\[
                         D_v
    =B_v-\frac1{N_v}A_v
\]
has degree zero and is a rational multiple of a principal divisor.  By
\cref{lem:p063-principal-residual-germ}, there is an ordinary smooth
harmonic germ \(h_v\), normalized by \(h_v(v)=G(v)\), whose graph slopes
are \(d_eG(v)\), whose slopes in the fresh directions \(p\in S_v\) are
\[
                         -a_{v,p},
             \qquad
                         a_{v,p}=\frac{m_{v,p}}{N_v}>0,
\]
and whose slopes in all remaining directions vanish.

We now choose one representative on one neighbourhood, rather than making
direction-by-direction choices.  Write it as a finite real linear
combination of logarithms of analytic units.  After one common shrink, each
unit has a rational reduction on \(C_v\); the union \(T_v\) of the supports
of their divisors is finite.  On every reduction fibre outside \(T_v\),
the absolute value of every unit is constant.  In a direction belonging to
\(T_v\) but not to the support of \(D_v\), the sum has outgoing slope zero;
after shrinking the finitely many such annular collars, its affine germ is
therefore constant.  This uses the reduction-fibre description of a simple
type-\(2\) neighbourhood in
\cite[\S4.29 and Lemma~5.12(3)]{BPR13Structure}.

Shrink once more along the finitely many graph and fresh directions.  On
this single neighbourhood, \(h_v\) agrees with \(G\circ\tau\)
simultaneously on every graph collar and on every zero-slope overlap: on a
graph collar the two functions have the same value and affine slope, and on
all remaining directions both are constant with value \(G(v)\).  The
avoidance of \(E_v\) ensures that each \(p\in S_v\)
determines an available open-disc component of \(Y\setminus\{v\}\),
disjoint from \(\Gamma\).

We cap the negative slope in each of these fresh directions.  Fix an even
nonnegative function \(\rho\in C^\infty_c((-1,1))\) of integral one and
put
\[
 \chi_\eta(r)=\int_{\mathbf R}|r-\eta s|\rho(s)\,\mathrm ds,
 \qquad
 m_\eta(s,t)=\frac{s+t+\chi_\eta(s-t)}2.
\]
Then \(m_\eta\) is smooth, convex, and nondecreasing in each variable,
\begin{equation}
\label{eq:p063-regularized-max-bounds}
 \max\{s,t\}\leq m_\eta(s,t)
       \leq\max\{s,t\}+\frac{\eta}{2},
\end{equation}
and it equals \(\max\{s,t\}\) if \(|s-t|\geq\eta\).  Indeed, Jensen's
inequality gives \(\chi_\eta(r)\geq|r|\), the triangle inequality gives
\(\chi_\eta(r)\leq|r|+\eta\), evenness gives zero first moment, and hence
\(\chi_\eta(r)=|r|\) for \(|r|\geq\eta\).  Also
\(|\chi_\eta'|\leq1\) and \(\chi_\eta''\geq0\).

Fix \(p\in S_v\).  On a sufficiently short punctured neighbourhood in its
open-disc component, choose an invertible annular coordinate \(T\) and a
constant \(c\) such that
\[
                         t=-\log|T|+c\geq0
\]
is distance from \(v\) on the skeleton.  Shrinking the whole annular tube,
not merely its skeletal arm, we have
\[
                         h_v=G(v)-a_{v,p}t.
\]
Choose \(r_{v,p}>0\) and \(\eta_{v,p}>0\) so that
\[
 \eta_{v,p}<a_{v,p}r_{v,p},
 \qquad
 r_{v,p}+\frac{\eta_{v,p}}{a_{v,p}}
\]
lies inside the arm on which this formula holds, and
\begin{equation}
\label{eq:p063-fresh-collar-error-choice}
 a_{v,p}r_{v,p}+\frac{\eta_{v,p}}2
 <\frac{\varepsilon}{3}.
\end{equation}
On that entire annular tube set
\begin{equation}
\label{eq:p063-fresh-collar-cap}
 q_{v,p}(t)
  =
 m_{\eta_{v,p}}\bigl(
     G(v)-a_{v,p}t,\,
     G(v)-a_{v,p}r_{v,p}
   \bigr).
\end{equation}
It equals \(h_v\) on
\[
               0\leq t<
               r_{v,p}-\frac{\eta_{v,p}}{a_{v,p}}
\]
and equals the constant \(G(v)-a_{v,p}r_{v,p}\) when
\[
               t>
               r_{v,p}+\frac{\eta_{v,p}}{a_{v,p}}.
\]
The same formula therefore assigns values on all residue discs attached to
the skeletal arm.  Extend it by that constant over the rest of the
open-disc component in the direction \(p\).  No return to the old constant
\(G(v)\) is made or needed.  By
\eqref{eq:p063-regularized-max-bounds},
\begin{equation}
\label{eq:p063-fresh-collar-error}
 \left|q_{v,p}-G(v)\right|
 \leq a_{v,p}r_{v,p}+\frac{\eta_{v,p}}2
 <\frac{\varepsilon}{3}.
\end{equation}

The function in \eqref{eq:p063-fresh-collar-cap} is ordinary smooth and
psh.  Indeed, \(t\) differs by a constant from \(-\log|T|\) for an
invertible analytic annular coordinate, and
\eqref{eq:p063-fresh-collar-cap} is a smooth convex function of \(t\) on
the whole annular tube.
Equivalently, it is a smooth convex coordinatewise nondecreasing
regularized maximum of the ordinary smooth harmonic functions \(h_v\)
and \(G(v)-a_{v,p}r_{v,p}\).  Near \(v\) every cap equals \(h_v\), while
at its far end it equals its constant extension on an open overlap.
Since \(S_v\) is finite, the modifications therefore glue around \(v\)
to an ordinary smooth psh function.

It remains to smooth the corners at interior type-\(3\) subdivision
vertices.  A smooth analytic curve has exactly two tangent directions at
such a point.  In an annular coordinate \(t\), with the vertex at \(0\),
write
\[
 G(t)=
 \begin{cases}
  G(v)+\alpha t,&t\leq0,\\
  G(v)+\beta t,&t\geq0.
 \end{cases}
\]
The outgoing-slope sum is \(\beta-\alpha\), so
\eqref{eq:p063-tree-positive-vertex-mass} gives \(\alpha<\beta\), and
\[
 G(t)=\max\{G(v)+\alpha t,G(v)+\beta t\}.
\]
Choose a collar on which the two affine formulas hold and which is disjoint
from every other replacement collar.  On that collar replace this maximum
by
\[
 m_\eta\bigl(G(v)+\alpha t,G(v)+\beta t\bigr).
\]
Choose \(\eta>0\) so small that
\(\eta/(\beta-\alpha)\) lies strictly inside this chosen collar.  The
replacement agrees with \(G\) whenever
\[
                         |t|>\frac{\eta}{\beta-\alpha},
\]
is ordinary smooth and psh, and differs from \(G\) by at most
\(\eta/2\).  Decrease \(\eta\), if necessary, so that also
\(\eta<2\varepsilon/3\).

There are only finitely many interior type-\(2\) and type-\(3\) vertices,
so all graph collars can be chosen pairwise disjoint.  At a type-\(2\)
vertex the central function is \(h_v\); on every graph or zero-slope arm
it equals \(G\circ\tau\) on an open overlap.  Every fresh direction is
off the graph, different fresh directions give disjoint open-disc
components, and the cap equals its constant continuation on an open
overlap.  The type-\(3\) replacement equals \(G\circ\tau\) on open
overlaps at both ends.  Away from these collars,
\(G\circ\tau\) is either affine in an ordinary annular coordinate or
locally constant.  Thus all local functions glue to a continuous
function \(g\) on \(Y\), and \(g|_{Y^\circ}\) is ordinary smooth and psh.
Boundary vertices need no smoothing because smoothness is asserted only
on \(Y^\circ\).

The local replacement supports are disjoint.  By
\eqref{eq:p063-fresh-collar-error} and the type-\(3\) error bound,
\[
                \lVert g-G\circ\tau\rVert_{L^\infty(Y)}
                    <\frac{\varepsilon}{3}.
\]
Together with \eqref{eq:p063-tree-graph-error}, this gives
\[
                         \lVert g-f\rVert_{L^\infty(Y)}
                         <\frac{2\varepsilon}{3}<\varepsilon,
\]
which proves \eqref{eq:p063-lisse-uniform-regularization}.
\end{proof}

\begin{proof}[Proof of
\cref{thm:p063-current-psh-ordinary-regularization}]
The equivalence between
\eqref{eq:p063-current-positivity} and Thuillier subharmonicity for
continuous functions is \cite[Corollary~4.8]{Wanner19}.
We prove that condition \(2\) implies condition \(3\).

Fix \(x\in W\).  By \cref{lem:p063-local-tree-affinoid}, choose a connected
strict affinoid tree domain \(Y\subset W\) with \(x\in Y^\circ\).  Choose
an open set \(\Omega\) with
\[
                         Y\subset\Omega\Subset W.
\]
Thuillier's approximation theorem
\cite[Theorem~3.4.2]{Thu05} gives a decreasing net
\((f_i)_i\) of lisse subharmonic functions on \(\Omega\), converging
pointwise to \(u\).  Its restriction to the compact space \(Y\) converges
uniformly.  Indeed, for \(\delta>0\), the increasing family of open sets
\[
                         U_i=\{y\in Y:f_i(y)-u(y)<\delta\}
\]
covers \(Y\).  A finite subcover and directedness provide one index \(i\)
with \(U_i=Y\).  This is the net form of Dini's argument; it uses the
continuity of \(u\) and of every \(f_i\).

Put \(b_n=4^{-n}\).  Recursively choose an increasing sequence of indices
\(i_n\) such that
\[
                         \lVert f_{i_n}-u\rVert_{L^\infty(Y)}
                         <\frac{b_n}{2}.
\]
At the \(n\)-th step, take a common upper bound of \(i_{n-1}\) and an
index supplied by uniform convergence.  Apply
\cref{prop:p063-tree-lisse-ordinary-regularization} to
\(f_{i_n}|_Y\), with error \(b_n/2\).  We obtain \(g_n\in C(Y)\), ordinary
smooth and psh on \(Y^\circ\), with
\begin{equation}
\label{eq:p063-raw-uniform-error}
                         \lVert g_n-u\rVert_{L^\infty(Y)}<b_n.
\end{equation}

The functions \(g_n\) need not be ordered.  Define
\begin{equation}
\label{eq:p063-tail-shift}
                         c_n
   =\sum_{j=n}^{\infty}(b_j+b_{j+1}),
   \qquad
                         u_n=g_n+c_n.
\end{equation}
Since \(c_n-c_{n+1}=b_n+b_{n+1}\),
\eqref{eq:p063-raw-uniform-error} gives, pointwise on \(Y\),
\[
 \begin{aligned}
 u_{n+1}-u_n
   &=
 g_{n+1}-g_n-(b_n+b_{n+1})\\
   &\leq
 |g_{n+1}-u|+|g_n-u|-(b_n+b_{n+1})
 \leq0.
 \end{aligned}
\]
Thus \(u_{n+1}\leq u_n\).  Moreover,
\[
                         \lVert u_n-u\rVert_\infty
                         \leq b_n+c_n\longrightarrow0.
\]
Adding constants preserves ordinary smoothness and psh.  This proves
\eqref{eq:p063-monotone-ordinary-approximation}.

Finally assume condition \(3\).  On \(Y^\circ\), each
\(\mathrm d'\mathrm d''[u_n]\) is a positive current.  Uniform convergence
gives \([u_n]\to[u]\) as currents: for every compactly supported smooth
test function \(\varphi\), the absolute value of
\[
 \bigl\langle [u_n-u],\mathrm d'\mathrm d''\varphi\bigr\rangle
\]
is at most \(\lVert u_n-u\rVert_\infty\) times the total variation of the
compactly supported smooth signed measure
\(\mathrm d'\mathrm d''\varphi\).  Hence
\[
                    \mathrm d'\mathrm d''[u_n]
                    \longrightarrow
                    \mathrm d'\mathrm d''[u]
\]
weakly.  Positivity is closed under weak convergence, so
\(\mathrm d'\mathrm d''[u]\geq0\) on \(Y^\circ\).  These interiors cover
\(W\), proving condition \(1\).

Condition \(3\) is exactly local membership, for continuous functions, in
the first decreasing-limit stage generated by ordinary smooth psh
functions.  The convergence obtained above is even uniform.
\end{proof}

\endgroup

\section{Reduction of cubical delta-metricity}
\label{sec:appN-cubical-reduction}

We prove \Cref{thm:p09p2-reduction}.  The theorem of the cube and the
symmetric--antisymmetric decomposition reduce the general question to
ample symmetric rigidified line bundles.

\begingroup
\small
% Full proof for the exact statement in
% p09p2_reduction_statement_exact_20260825.tex.

\paragraph{Imported inputs.}
The proof uses the following statements with exactly the hypotheses recorded
here.
\begin{enumerate}
\item Gubler--K\"unnemann work over an algebraically closed, complete,
nontrivially valued non-Archimedean field.  Their Proposition~8.11
identifies formal and piecewise linear metrics; their Proposition~8.13
identifies algebraic and formal metrics on line bundles over proper
varieties; their Definition~9.9 and \S9.13 define ordinary delta-metrics in
algebraic tropical charts; their \S9.14 gives isometry invariance, tensor
and dual closure, and descent from a positive tensor power; and their
Remark~9.16 makes piecewise linear, formal, and algebraic metrics
delta-metrics.  Their Example~8.16 states that the normalized canonical
metric of an algebraically trivial line bundle has an algebraic positive
tensor power.

\item The algebraicity statement used in
\cite[Example~8.16]{GK17} is \cite[Example~3.7]{Gub10CM}.  It is used only
in the form quoted by Gubler--K\"unnemann.

\item Chambert--Loir--Ducros,
\cite[Lemmas~19.1.4 and 19.1.6]{CLD25}, give compatibility of cubical
metrics with tensor products, duals, and homomorphism pullbacks, existence
and uniqueness up to a positive scalar, uniqueness after normalizing a
rigidification, and the exponent-one characterization for an
antisymmetric line bundle.

\item Milne, \cite[Proposition~8.4 and Remark~8.5]{milneAV}, gives the
characterization of \(\operatorname{Pic}^0(A)\) used below and inversion as
duality on \(\operatorname{Pic}^0(A)\); his
\cite[Remark~10.19]{milneAV} states that inversion acts trivially on
\(\operatorname{NS}(A)\).

\item The Stacks Project, \cite[Lemma~39.9.2]{stacks-project}, gives an
ample line bundle on an abelian variety.  Its Lemma~28.27.5
(Tag~0890) preserves ampleness under the tensor product used below, and
its Lemma~29.38.7 (Tag~0892) gives
\(S\otimes H^{\otimes m}\) ample for all sufficiently large \(m\) when
\(H\) is ample.
\end{enumerate}

\paragraph{Metric recognition and closure rules.}
The proof uses the following meanings and no identifications between them
without a cited recognition result.  An algebraic metric on a line bundle
over a proper variety is induced by a flat proper finite-type
\(K^\circ\)-model and a line bundle on that model.  A formal metric is
induced by an admissible formal \(K^\circ\)-model.  On a proper variety,
algebraic and formal metrics coincide by
\cite[Proposition~8.13]{GK17}; formal and piecewise linear metrics coincide
by \cite[Proposition~8.11]{GK17}.  A piecewise linear metric is an ordinary
delta-metric by \cite[\S9.13 and Remark~9.16]{GK17}.

For comparison with the local definition, a piecewise smooth weight
\(\phi\) in an ordinary tropical chart is a delta-weight when its facewise
Hessian
\[
                 \mathrm d'_{\mathrm P}\mathrm d''_{\mathrm P}\phi
\]
is induced by a codimension-zero generalized delta-preform on the ambient
tropical vector space.  This is the condition in
\cite[Definition~9.9 and \S9.13]{GK17}.  In particular, a constant real
weight is smooth and is a delta-weight, whether or not the constant lies in
the value group of \(K\).

We shall repeatedly use the exact closure assertions of
\cite[\S9.14]{GK17}: ordinary delta-metrics are invariant under isometry,
are closed under tensor product and duality, and satisfy
\begin{equation}
\label{eq:p09p2-root-rule}
 \overline M^{\otimes N}\text{ is an ordinary delta-metric for some }N\geq1
 \quad\Longrightarrow\quad
 \overline M\text{ is an ordinary delta-metric}.
\end{equation}
Thus \eqref{eq:p09p2-root-rule} is a theorem about the already given metric
on \(M\); it neither asserts nor requires the existence of an \(N\)-th root
of an unrelated line bundle or metric.

\begin{proof}
\cref{itm:p09p2-pic0}
Choose a rigidification \(\varepsilon\in L(0)\setminus\{0\}\).  Multiplying
the given metric by a positive constant, if necessary, gives a cubical
metric \(\|\cdot\|_0\) for which \(\|\varepsilon\|_0=1\).  Constant
rescaling preserves the cubical condition: the total exponent with which a
constant enters the metric on \eqref{eq:p09p2-d3} is
\[
        \sum_{I\subseteq\{1,2,3\}}(-1)^{\# I+1}=0.
\]

Since \(L\in\operatorname{Pic}^0(A)\), one has
\(\iota^*L\simeq L^\vee\).  Indeed, the characterization of
\(\operatorname{Pic}^0(A)\) in \cite[Proposition~8.4]{milneAV} gives
\(m^*L\simeq p_1^*L\otimes p_2^*L\) on \(A\times A\); pulling this
isomorphism back by \(x\mapsto(-x,x)\) yields
\(\iota^*L\otimes L\simeq\mathcal O_A\).  The normalized cubical metric
\(\|\cdot\|_0\) is therefore the normalized canonical metric used for an
algebraically trivial line bundle in \cite[Example~8.16]{GK17}; equivalently,
it is characterized by the isometric relation
\([n]^*\overline L_0\simeq\overline L_0^{\otimes n}\) for \(n\geq2\), as in
\cite[Lemmas~19.1.4 and 19.1.6]{CLD25}.

By \cite[Example~8.16]{GK17}, whose algebraicity input is
\cite[Example~3.7]{Gub10CM}, there is an integer \(N\geq1\) for which
\(\overline L_0^{\otimes N}\) is an algebraic metric.  Because \(A\) is
proper, the recognition chain recorded above gives
\[
 \overline L_0^{\otimes N}\text{ algebraic}
 \Longrightarrow
 \overline L_0^{\otimes N}\text{ formal}
 \Longrightarrow
 \overline L_0^{\otimes N}\text{ piecewise linear}
 \Longrightarrow
 \overline L_0^{\otimes N}\text{ an ordinary delta-metric}.
\]
Applying \eqref{eq:p09p2-root-rule} proves that \(\overline L_0\) is an
ordinary delta-metric.

Return to the original cubical metric.  It has the form
\(\overline L=\overline L_0\otimes\overline{\mathcal O}_{c}\), where
\(\overline{\mathcal O}_{c}\) is the trivial line bundle with a constant
positive norm.  Its local weight is constant, hence smooth and an ordinary
delta-weight.  Tensor closure now proves \cref{itm:p09p2-pic0} for the
original, arbitrarily normalized cubical metric.

\cref{itm:p09p2-split}
The operations in \eqref{eq:p09p2-plus-minus} preserve cubical metrics.
This follows directly from \eqref{eq:p09p2-d3}: \(D_3\) commutes with tensor
products and duals, and for a homomorphism \(u:A\to A\) one has
\(D_3(u^*L)=(u^3)^*D_3(L)\), compatibly with the distinguished
trivializations.  This is also the tensor, dual, and pullback compatibility
of \cite[Lemmas~19.1.4 and 19.1.6]{CLD25}.

Because \(\iota^2=\operatorname{id}_A\), the definitions give
\[
 \iota^*L_+
   =\iota^*L\otimes L\simeq L_+,
 \qquad
 \iota^*L_-
   =\iota^*L\otimes L^\vee\simeq L_-^\vee.
\]
Moreover, \(\iota^*\) acts as the identity on
\(\operatorname{NS}(A)=\operatorname{Pic}(A)/\operatorname{Pic}^0(A)\),
by \cite[Remark~10.19]{milneAV}.  Hence
\[
 [L_-]_{\operatorname{NS}}
   =[L]_{\operatorname{NS}}-[\iota^*L]_{\operatorname{NS}}=0,
\]
so \(L_-\in\operatorname{Pic}^0(A)\).  Finally, evaluation for the pullback
metric gives the isometry
\[
 \begin{aligned}
 \overline L_+\otimes\overline L_-
  &=(\overline L\otimes\iota^*\overline L)
    \otimes(\overline L\otimes(\iota^*\overline L)^\vee)\\
  &\simeq \overline L^{\otimes2},
 \end{aligned}
\]
which is \eqref{eq:p09p2-square-split}.  The rigidifications on
\(L_+\), \(L_-\), and \(L^{\otimes2}\) are the induced ones, and there is no
normalization scalar in this isometry because all four metrics are the
directly induced tensor, pullback, and dual metrics from
\eqref{eq:p09p2-plus-minus}.  In particular, when
\(\|\varepsilon\|_L=1\), the induced rigidifications on \(L_+\), \(L_-\),
and \(L^{\otimes2}\) also have norm one and the displayed isometry is an
isometry in the normalized rigidified category.

By \cref{itm:p09p2-pic0}, \(\overline L_-\) is an ordinary delta-metric.  If
\(\overline L_+\) is an ordinary delta-metric, tensor closure and
\eqref{eq:p09p2-square-split} make \(\overline L^{\otimes2}\) an ordinary
delta-metric, and \eqref{eq:p09p2-root-rule} makes \(\overline L\) one.  In
the reverse direction, if \(\overline L\) is an ordinary delta-metric, then
\[
       \overline L_+
       \simeq
       \overline L^{\otimes2}\otimes\overline L_-^\vee
\]
is an ordinary delta-metric by tensor and dual closure.  This proves
\eqref{eq:p09p2-delta-equivalence}.

\cref{itm:p09p2-ample-reduction}
The implication
\cref{itm:p09p2-all}\(\Rightarrow\)\cref{itm:p09p2-ample-symmetric} is
immediate.  Assume \cref{itm:p09p2-ample-symmetric}.  We first prove that
every cubical metric \(\overline S\) on every symmetric line bundle \(S\)
is an ordinary delta-metric.

An abelian variety is projective by
\cite[Lemma~39.9.2]{stacks-project}, so choose an ample line bundle \(H_0\)
on \(A\), choose a rigidification, and give it its normalized cubical metric
\(\overline H_0\), whose existence and uniqueness are
\cite[Lemmas~19.1.4 and 19.1.6]{CLD25}.  Set
\[
                 \overline H
        :=\overline H_0\otimes\iota^*\overline H_0.
\]
The preceding compatibility shows that \(\overline H\) is cubical, and
\(\iota^*H\simeq H\).  Pullback by the automorphism \(\iota\) preserves
ampleness, and the tensor product of ample line bundles is ample by
\cite[Lemma~28.27.5, Tag~0890]{stacks-project}; hence \(H\) is ample.

For all sufficiently large integers \(m\), the line bundle
\begin{equation}
\label{eq:p09p2-ample-twist}
                         S_m:=S\otimes H^{\otimes m}
\end{equation}
is ample.  This is
\cite[Lemma~29.38.7, Tag~0892]{stacks-project}, applied to the identity
morphism of \(A\): every invertible sheaf is relatively ample for the
identity morphism, and \(H\) is ample on the target.  Both \(H\) and \(S_m\)
are symmetric.  Give \(S_m\) the induced rigidification and cubical metric
\(\overline S_m=\overline S\otimes\overline H^{\otimes m}\).  The assumed
statement \cref{itm:p09p2-ample-symmetric} applies to \(\overline H\) and
\(\overline S_m\), so both are ordinary delta-metrics.  Therefore
\[
          \overline S
          \simeq
          \overline S_m\otimes
          (\overline H^{\otimes m})^\vee
\]
is an ordinary delta-metric by tensor and dual closure.

Now let \(\overline L\) be any cubically metrized line bundle on \(A\).
The line bundle \(L_+\) of \eqref{eq:p09p2-plus-minus} is symmetric, so the
preceding paragraph makes \(\overline L_+\) an ordinary delta-metric.
Equation \eqref{eq:p09p2-delta-equivalence} then makes \(\overline L\) an
ordinary delta-metric.  Thus
\cref{itm:p09p2-ample-symmetric}\(\Rightarrow\)\cref{itm:p09p2-all}, and
the proof is complete.
\end{proof}

\endgroup

\section{The product--Gauss criterion}
\label{sec:appN-product-gauss}

We prove \Cref{thm:p09p2-product-gauss-candidate}.  The proof combines
algebraic action pullback, relative Gauss points, common-cell lineality, and
specialization to characterize ordinary smoothness of the Raynaud
tropicalization.

\begingroup
\small
% Exact candidate proof for source and production-Rethlas audit.
% This file is isolated from the live book. It makes no novelty or priority
% claim and must not be promoted before a final/correct production receipt.

\begin{proof}
Put \(\Gamma=-\log|K^\times|\).  We first record the three chart facts
used below.  A finite integral piecewise-linear tropicalization on a good
strict analytic space is locally covered by one smooth analytic moment map
and a piecewise integral-affine map; see
\cite[Proposition~8.3 and \S8.4]{GJR26}.  On an open subset of an algebraic
analytification, \cite[Proposition~7.2]{Gub16} replaces such a finite
analytic moment, after shrinking, by an algebraic moment with the same
tropicalization.  It is equality of the tropicalizations, rather than of
the analytic functions themselves, that will be used.  By factoring the
algebraic moment through the canonical moment map of a common very affine
open, as in \cite[\S4.12]{Gub16}, one obtains an ordinary algebraic
tropical chart dominating any finitely many prescribed charts.

The third fact is the full-ambient condition in the definition of an
ordinary delta-metric.  In an ordinary delta-chart, the facewise Hessian of
the weight is the restriction of a \((1,1)\)-superform on the full ambient
vector space with piecewise-smooth coefficients
\cite[Definition~9.9(ii) and \S9.13]{GK17}.  A piecewise-smooth coefficient
is one single-valued continuous function across the cells
\cite[Definition~1.6]{GK17}; equivalently, the ambient superform is a
codimension-zero delta-preform
\cite[Remark~1.7, Example~2.10, and Remark~9.10]{GK17}.  Thus its
coefficient tensor has one value at a wall, not one unrelated value for
each incident face.

We next establish the specialization step that will return from a product
Gauss point to the original variety.  Let \(Y\) be an algebraic
\(K\)-variety, let \(y\in Y^{\mathrm{an}}\), and write
\(H_y=\mathcal H(y)\).  For \(w\in N_{\mathbf R}\), let

\[
             \zeta_{w,y}\in T_{H_y}^{\mathrm{an}}
                    \subset (T\times Y)^{\mathrm{an}}
\]

be the relative Gauss point characterized by

\begin{equation}
\label{eq:p09p2-relative-gauss-point}
 \left|\sum_{m\in M}a_m\chi^m\right|_{\zeta_{w,y}}
   =\max_m |a_m|\exp(-\langle m,w\rangle),
 \qquad a_m\in H_y.
\end{equation}

This is a point of the analytic fibre product, not a naive ordered pair of
Berkovich points.  The Shilov-section construction of
\cite[\S\S12.1.10--12.1.11]{CLD25} also shows that

\[
             (w,y)\longmapsto\zeta_{w,y}
\]

is continuous.

We claim that, if \(w\in N\otimes_{\mathbf Z}\Gamma\), every ordinary
analytic neighbourhood \(W\) of \(\zeta_{w,y}\) contains \(i_t(y)\) for
some \(t\in T(K)\) satisfying \(\operatorname{trop}_T(t)=w\), where

\[
                       i_t:Y\longrightarrow T\times Y,
                       \qquad z\longmapsto(t,z)
\]

is the constant algebraic section.  Choose a basis of \(M\).  Since
\(w\in N\otimes\Gamma\), scaling the corresponding torus coordinates by
elements of \(K^\times\) reduces the assertion to the unit Gauss point of

\[
 \mathcal T_y=\mathcal M\bigl(
 H_y\langle Z_1^{\pm1},\ldots,Z_r^{\pm1}\rangle\bigr).
\]

The intersection \(W\cap\mathcal T_y\) contains a basic ordinary
neighbourhood of the Gauss point controlled by the absolute values of
finitely many analytic functions.  After normalizing each nonzero function
by a coefficient attaining its Gauss norm, its initial reduction is a
nonzero Laurent polynomial over \(\widetilde H_y\).  It is therefore enough
to choose one element of \((\widetilde K^\times)^r\) outside the zero loci
of finitely many nonzero Laurent polynomials with coefficients in
\(\widetilde H_y\).

For completeness, \((\widetilde K^\times)^r\) is Zariski dense in
\((\widetilde H_y^\times)^r\).  More generally, if \(k\subset L\) are
fields with \(k\) infinite, a nonzero Laurent polynomial over \(L\) cannot
vanish on all of \((k^\times)^r\).  This follows by induction on \(r\):
choose the first \(r-1\) entries so that one nonzero coefficient polynomial
in the last variable remains nonzero, and then avoid the finitely many
roots of the resulting one-variable polynomial.  Here the isometric
embedding \(K\hookrightarrow H_y\) induces
\(\widetilde K\hookrightarrow\widetilde H_y\), while \(\widetilde K\) is
algebraically closed and hence infinite.  Lifting the chosen residue tuple
to \((K^{\circ\times})^r\), the ultrametric inequality gives equality with
the Gauss norm for every one of the finitely many functions.  Multiplying
back the scaling coordinates produces the required \(t\in T(K)\).  This
proves the claim without invoking rigid-point density.

Consequently, let \(u\) be a real-valued function germ at \(y\) and let
\(m\in M_{\mathbf R}\).  If

\begin{equation}
\label{eq:p09p2-gauss-specialization}
 \operatorname{pr}_Y^*u
   +\langle m,\operatorname{trop}_T\circ\operatorname{pr}_T\rangle
\end{equation}

is CLD-smooth at \(\zeta_{w,y}\), choose a neighbourhood \(W\) on which
it is smooth and then choose \(t\) as above with \(i_t(y)\in W\).  The set
\(i_t^{-1}(W)\) is an ordinary neighbourhood of \(y\).  Pulling
\eqref{eq:p09p2-gauss-specialization} back by \(i_t\) shows there that

\[
 u+\langle m,\operatorname{trop}_T(t)\rangle
\]

is smooth.  The second term is constant, so \(u\) is smooth at \(y\).

Assume now that the metric with weight \(F\circ h\) is an ordinary
delta-metric.  Put

\[
 \rho=\operatorname{trop}_T\circ\operatorname{pr}_T
           +h\circ\operatorname{pr}_X
       :(T\times X)^{\mathrm{an}}\longrightarrow N_{\mathbf R}.
\]

The equivariance identity \eqref{eq:p09p2-product-equivariance} gives

\begin{equation}
\label{eq:p09p2-action-pullback-weight}
                         F\circ\rho
                    =(F\circ h)\circ\mu^{\mathrm{an}}.
\end{equation}

Algebraic pullback preserves ordinary delta-metrics by
\cite[Proposition~9.15]{GK17}.  Hence \(F\circ\rho\) is an ordinary
delta-weight on \(T\times X\).

Fix \(x\in X^{\mathrm{an}}\), choose provisionally
\(w_0\in N\otimes\Gamma\), and take an ordinary delta-chart at
\(\zeta_{w_0,x}\).  Cover all coordinates of the locally integral PL map
\(\rho\) by one analytic moment, algebraize that finite moment by
\cite[Proposition~7.2]{Gub16}, adjoin a basis of characters of the first
\(T\)-factor, and pass to one common ordinary algebraic refinement.  Since
the line bundle is trivial, we use its global frame \(1\); if the original
delta-presentation used another algebraic frame, adjoining their algebraic
unit ratio changes the weight only by an affine tropical coordinate and
does not change its facewise Hessian.  After shrinking, we obtain

\[
 \Phi:V\longrightarrow P\subset\mathbf R^s,
 \qquad
 H:P\longrightarrow N_{\mathbf R},
\]

where \(\Phi\) is the tropicalization in an ordinary algebraic chart,
\(H\) is piecewise integral affine, and

\begin{equation}
\label{eq:p09p2-common-chart}
 \rho|_V=H\circ\Phi,
 \qquad
 \mathrm d'_{\mathrm P}\mathrm d''_{\mathrm P}(F\circ H)
          =\gamma|_P
\end{equation}

for a \((1,1)\)-superform \(\gamma\) on \(\mathbf R^s\) with
piecewise-smooth coefficients.  The old ambient superform has been pulled
back along the integral-affine projection from the refined chart.

Fix common locally finite complexes of definition for \(H\), the tropical
support, and every coefficient function of \(\gamma\), and restrict to a
small target neighbourhood meeting only finitely many cells.  All chart
coordinates, factorizations, and walls have now been fixed; only now do we
choose the final Gauss weight.

For \(w'\) sufficiently close to \(w_0\), continuity of the Shilov section
keeps \(\zeta_{w',x}\) in \(V\).  Set

\[
                         z(w')=\Phi(\zeta_{w',x}).
\]

Each algebraic unit defining the final chart restricts on the fibre over
\(x\) to an element of \(\mathcal H(x)(M)^\times\), hence to a quotient of
two finite Laurent polynomials.  By \eqref{eq:p09p2-relative-gauss-point},
its logarithmic norm is the difference of the maxima of finitely many
integral-affine functions of \(w'\).  Thus \(z\) is finite piecewise
integral affine on a sufficiently small neighbourhood of \(w_0\).

Refine its finite source subdivision by the inverse images of all the
already fixed target cells.  Since \(K\) is algebraically closed,
\(\Gamma\) is nonzero and divisible; it is therefore dense in \(\mathbf R\).
We may choose

\[
                         w\in N\otimes_{\mathbf Z}\Gamma
\]

arbitrarily close to \(w_0\), in the relative interior of a
full-dimensional cell of the source refinement.  Put \(z_0=z(w)\).  There
are an open neighbourhood \(D\) of \(w\), an injective linear map

\begin{equation}
\label{eq:p09p2-gauss-slope-map}
                         a:N_{\mathbf R}\hookrightarrow\mathbf R^s,
\end{equation}

and a cell \(C\) of the common target refinement such that

\begin{equation}
\label{eq:p09p2-gauss-common-cell}
 z(w+u)=z_0+a(u),
 \qquad
 z_0\in\operatorname{relint}(C),
 \qquad
 z(D)\subset C.
\end{equation}

The character coordinates of the first torus factor make \(a\) injective.
Since \(D\) is full dimensional, \eqref{eq:p09p2-gauss-common-cell} gives

\begin{equation}
\label{eq:p09p2-gauss-lineality}
                         a(N_{\mathbf R})\subset L_C.
\end{equation}

Moreover, for \(u\) sufficiently small,

\[
 H(z(w+u))=\rho(\zeta_{w+u,x})=w+u+h(x).
\]

Differentiation on \(C\) gives \(D_CH\circ a=\operatorname{id}\).  If
\(\sigma\) is any cell in the star of \(C\), then \(C\) is a face of
\(\sigma\), the affine restrictions \(H|_C\) and \(H|_\sigma\) agree on
\(C\), and \(a(N_{\mathbf R})\subset L_C\subset L_\sigma\).  Hence

\begin{equation}
\label{eq:p09p2-H-lineality}
              D_\sigma H(a(v))=v
 \qquad
 (v\in N_{\mathbf R},\ \sigma\in\operatorname{Star}(C)).
\end{equation}

We now use the full-ambient tensor in
\eqref{eq:p09p2-common-chart}.  Let \(G_0\) be the coefficient bilinear
form of \(\gamma\) at \(z_0\).  On the relative interior of a cell
\(\sigma\) in the star of \(C\), the facewise Hessian of \(F\circ H\) is

\[
              b(D_\sigma H\,\mathord\cdot,
                D_\sigma H\,\mathord\cdot).
\]

The coefficient functions of \(\gamma\) have one common value at
\(z_0\).  Taking the limit to \(z_0\) through the relative interior of
\(\sigma\) therefore yields

\begin{equation}
\label{eq:p09p2-wall-tensor}
 \left.G_0\right|_{L_\sigma\times L_\sigma}
 =b(D_\sigma H\,\mathord\cdot,D_\sigma H\,\mathord\cdot).
\end{equation}

Fix \(m\in M\), choose \(v=(b^\sharp)^{-1}(m)\), where

\[
 b^\sharp:N_{\mathbf R}\longrightarrow M_{\mathbf R},
 \qquad v\longmapsto b(v,\mathord\cdot),
\]

and put

\[
                         \alpha_m=G_0(a(v),\mathord\cdot).
\]

By \eqref{eq:p09p2-gauss-lineality}, \eqref{eq:p09p2-H-lineality}, and
\eqref{eq:p09p2-wall-tensor},

\[
 \left.\alpha_m\right|_{L_\sigma}
   =b(v,D_\sigma H\,\mathord\cdot)
   =\left.\mathrm d\langle m,H\rangle\right|_\sigma.
\]

The function \(\langle m,H\rangle\) is continuous and affine on every
cell, and every cell in the star contains \(C\) as a face.  Its value at
\(z_0\) and its derivative on every incident cell therefore agree with one
ambient affine function.  On the whole star germ,

\begin{equation}
\label{eq:p09p2-affine-coordinate}
 \langle m,H(z)\rangle
  =\langle m,H(z_0)\rangle+\alpha_m(z-z_0).
\end{equation}

It follows that \(\langle m,\rho\rangle\) is CLD-smooth at
\(\zeta_{w,x}\).  Since

\[
 \langle m,\rho\rangle
 =\langle m,\operatorname{trop}_T\circ\operatorname{pr}_T\rangle
    +\operatorname{pr}_X^*\langle m,h\rangle,
\]

the specialization argument applied to
\(u=\langle m,h\rangle\) shows that this coordinate is smooth at \(x\).
The point \(x\) and the lattice coordinate \(m\) were arbitrary, so every
coordinate of \(h\) is CLD-smooth.

Conversely, suppose that the coordinates of \(h\) in a basis of \(M\) are
CLD-smooth.  Near any point, concatenate their finitely many analytic
moment presentations into one analytic moment.  The coordinate vector
\(h\), and hence \(F\circ h\), is the pullback of a smooth function on its
tropical target.  Apply \cite[Proposition~7.2]{Gub16} once to replace the
finite analytic moment by an algebraic moment with the same
tropicalization, and factor it through the canonical moment of a very
affine open.  Thus \(F\circ h\) is smooth in the ordinary algebraic-chart
sense.  Smooth metrics are ordinary delta-metrics by
\cite[Remark~9.16]{GK17}.  Real linear combinations show that testing one
basis of \(M\) is equivalent to testing all coordinates.  This proves the
first assertion.

Before applying it to Raynaud uniformization, we verify one locality fact
not contained in \cite[Proposition~9.15]{GK17}.  Let

\[
                 \iota:V_X\xrightarrow{\sim}V_Y
\]

be an analytic isomorphism between analytic open subsets of algebraic
\(K\)-varieties.  Suppose algebraic line bundles on the two varieties are
identified over these opens by an analytic isomorphism, and transport a
metric through this identification.  Then the metric is ordinary delta on
one side if and only if it is ordinary delta on the other.

Indeed, pull the finitely many algebraic moment coordinates of a
delta-presentation on \(V_Y\) back by \(\iota\), and adjoin the invertible
analytic function relating the pulled-back algebraic frame to an algebraic
frame on the source.  Package this finite family as one analytic moment.
By \cite[Proposition~7.2]{Gub16}, after shrinking it has the same
tropicalization as an algebraic moment.  Factoring through the canonical
moment map gives an ordinary algebraic chart.  On its target, the old
weight is composed with an integral-affine projection and is changed only
by the last affine coordinate coming from the frame ratio; the old ambient
codimension-zero superform pulls back along that projection.  This is an
ordinary delta-presentation on the source.  Applying the same argument to
\(\iota^{-1}\) proves the equivalence.  The assertion is only being made
for analytic opens of algebraic varieties and algebraic line bundles.

We now consider the split Raynaud data in the statement.  The Raynaud
extension \(E\) is algebraic and its \(T\)-action is algebraic by
\cite[\S16]{GJR26}; compare \cite[Remark~B.6.6]{CLD25}.  In the split
situation, the quotient

\[
                       p:E^{\mathrm{an}}\longrightarrow A^{\mathrm{an}}
\]

is an analytic covering by \(\Lambda\), hence locally an analytic
isomorphism; see \cite[\S19.2.1 and \S\S B.5.12--B.5.13]{CLD25}.  The
locality fact just proved applies to the algebraic line bundle on \(E\)
in the following structure formula.

By \cite[Proposition~19.2.5]{CLD25}, there are a line bundle \(M_L\) on
\(B\), a formal metric on \(q^*M_L\), an isomorphism of its analytification
with \(p^*L\), and a quadratic function \(F_L\) on \(N_{\mathbf R}\) such
that the pullback of the normalized cubical metric is the product of that
formal metric and the metric with weight

\begin{equation}
\label{eq:p09p2-raynaud-weight}
                              F_L\circ\tau.
\end{equation}

Let \(Q_L\) denote the homogeneous quadratic form of Raynaud's
polarization data, and let \(b_L\) be its associated symmetric bilinear
form, normalized by \(Q_L(w)=b_L(w,w)\).  The same proposition says that
the homogeneous degree-two part of \(F_L\) is \(\frac12Q_L\); hence

\begin{equation}
\label{eq:p09p2-raynaud-quadratic}
             F_L(w)=\frac12b_L(w,w)+\ell_L(w)+c_L.
\end{equation}

The canonical tropicalization

\[
                       \tau:E^{\mathrm{an}}\longrightarrow N_{\mathbf R}
\]

is locally integral PL and satisfies

\[
              \tau(t\mathbin{\cdot}e)
                 =\operatorname{trop}_T(t)+\tau(e).
\]

Since \(L\) is ample, \cite[Theorem~B.6.4]{CLD25} says that \(Q_L\) is
positive definite; thus \(b_L\) is nondegenerate.  Formal metrics are
piecewise linear by \cite[Proposition~8.11]{GK17}, hence ordinary
delta-metrics by \cite[Remark~9.16]{GK17}.  Ordinary delta-metricity is
local, invariant under isometry, and closed under tensor products and duals
by \cite[\S9.14]{GK17}.  The analytic-local equivalence above and
\eqref{eq:p09p2-raynaud-weight} therefore give

\[
 \overline L_{\mathrm{cub}}\text{ is ordinary delta on }A^{\mathrm{an}}
 \quad\Longleftrightarrow\quad
 F_L\circ\tau\text{ is ordinary delta on }E^{\mathrm{an}}.
\]

Apply the first assertion with \(X=E\), \(h=\tau\), and \(b=b_L\).
It follows that the two equivalent conditions in the preceding display
hold if and only if every coordinate of \(\tau\) is CLD-smooth.  Finally,
\cite[Remark~19.2.3]{CLD25} identifies simultaneous smoothness of these
coordinates with isotriviality of the residual extension of \(B_s\) by
\(T_{\widetilde K}\).  Every cubical metric on a fixed line bundle
is a positive constant multiple of the normalized one
\cite[Lemma~19.1.6(b)]{CLD25}; multiplication by a constant changes the
weight by a
smooth constant and preserves ordinary delta-metricity.  This proves the
claimed equivalence and its non-isotrivial mixed-reduction consequence.
\end{proof}

\endgroup

\section{How to use the appendix}
\label{sec:appN-use}

The curve results are local analytic statements: they compare positivity of
\(\dprimeddprime[u]\), Thuillier subharmonicity, and ordinary smooth
approximation.  The cubical results are global statements about a metrized
line bundle on an abelian variety.  In particular, the product--Gauss
criterion does not assert that every harmonic tropicalization is ordinary
smooth; it identifies the precise residual isotriviality condition in the
Raynaud setting.
